\subsection{Firma documental individual con clave externa}
\label{sec:pruebas-firma-individual}

La verificación local del 10 de octubre de 2026 comprobó el vínculo de una
cuenta Litigator con un certificado público y la firma de una versión PDF
exacta mediante dos pruebas externas. El registro canónico identifica la
cuenta y su certificado; el consentimiento documental incorpora expediente,
documento, versión, operación, confianza y resumen de la firma del PDF.
La identidad de la cuenta no se deduce del nombre común del certificado.
Los resultados se separan por frontera en la
\cref{tab:matriz-firma-individual}; no constituyen una nueva medición global
de cobertura ni reemplazan los cortes históricos de este capítulo.

\subsubsection{Valores, emisión y pruebas criptográficas}

Doce pruebas de dominio aprobaron los vectores literales de 142 bytes del
vínculo \texttt{DOCBIND1} y de 258 bytes del consentimiento \texttt{DOCSIGN1},
su reconstrucción estricta y el rechazo de identidades, contadores y propósitos
alterados. Cinco pruebas de emisión mediante la CLI aprobaron con OpenSSL:
la hoja conserva la clave pública del CSR y el perfil interno de firma,
sin generar una clave privada del titular dentro de la CA. Se rechazaron
solicitudes alteradas, excesivas o privadas, parámetros RSA no admitidos,
sobrescritura de archivos y sustitución de la clave pública emitida.

Ocho casos criptográficos comprobaron las dos firmas RSA de 3\,072 bits,
certificado, confianza exacta, vigencia y revocación. La firma del
consentimiento no sustituye la comprobación de la firma separada del PDF.
Otro documento, versión, cuenta o propósito no puede reutilizar esa intención.
Once casos del servicio de vinculación y catorce del servicio documental
comprobaron reautenticación del principal completo, recibos exactos, cambios
de cuenta o confianza, retiro, respuesta perdida y ausencia de una segunda
escritura en la repetición autorizada.

Ocho pruebas del procesador comprobaron la conservación del material personal,
la TSA configurada, el rechazo de sellos de tiempo inválidos y la evaluación
del certificado en el instante de aceptación. Seis casos de evidencia
histórica rechazaron cambios de material, contexto o formato. La exportación
aprobó once pruebas, incluidas las cuatro de la extensión individual; el
repositorio de archivos aprobó diez, conservando la lectura de sellos anteriores.
Siete pruebas HTTP con puertos controlados verificaron rutas, límites,
identificadores, errores y la distinción explícita de evidencia del servidor.
Estos ensayos controlados no acreditan por sí solos el recorrido de servicios
reales descrito a continuación.

\subsubsection{Atomicidad, concurrencia y restauración nativa}

Con PostgreSQL~18.6 aislado y roles operativos restringidos aprobaron diecisiete
casos distintos: siete de vínculos, cinco del registro documental y cinco de
invariantes. El conteo reúne la campaña inicial y las repeticiones focales de
los casos corregidos, sin contar dos veces una identidad de prueba. Se
comprobaron permisos actuales, cambios de cuenta o confianza, huellas retenidas
por su cuenta original, un único vínculo activo, retiro terminal, concurrencia
y rollback de documento, recibo y auditoría ante un fallo de escritura.

Una evidencia personal sin recibo y una divergencia entre ambos se rechazaron
al confirmar la transacción. Desactivar cualquiera de los dos disparadores
diferidos impidió reabrir el adaptador. El volcado y la restauración lógica
conservaron la evidencia y sus eventos después de retirar el vínculo. Una
regresión reprodujo además la sustitución de la CRL capturada por una sucesora
válida: el vínculo se había registrado con la primera revisión y el documento
firmado con la segunda. El inventario corregido exige que la confianza
embebida coincida con la revisión exacta del recibo, y rechazó la sustitución
por la tercera. La concurrencia nativa conservó un solo recibo. En aplicación
se comprobó además que prevalecía la primera evidencia aceptada aunque ambas
solicitudes prepararan sellos de tiempo distintos.

Los primeros intentos de migración y la prueba que suponía un nombre automático
de restricción no se contabilizan como aprobaciones. Sus correcciones se
verificaron antes del resultado anterior. La admisión de formato y la TSA de
estas pruebas nativas utilizan puertos controlados; su finalidad es comprobar
el contrato transaccional con PostgreSQL real.

\subsubsection{Aceptación HTTP, evidencia pública y continuidad}

El recorrido focal HTTP aprobado utilizó PostgreSQL~18.6, Valkey~8.1.10 como
almacén compatible con Redis, ingreso real por contraseña y MFA, CA interna,
OpenSSL y TSA local real. La clave sintética se generó en un directorio del
cliente separado de la CA; el servidor recibió certificado y pruebas públicas.
La campaña ejecutó los pasos de
\rutarepo{scripts/api-individual-signing-demo.sh} y comprobó independientemente
los bytes canónicos mediante su auxiliar Python.

El litigante vinculó su certificado, descargó el contenido exacto de la versión
uno y aportó sus dos firmas. La versión dos del mismo documento permaneció
sin sellar. Se descartó deliberadamente la respuesta de confirmación y se
recuperó el recibo por su operación; la repetición devolvió el mismo resultado
sin otro registro de firma. OpenSSL verificó las pruebas públicas del ZIP,
la cadena y CRL capturadas y el token RFC~3161 sobre el PDF.

Después se publicó una CRL que revocaba la hoja. La verificación mantuvo la
validez histórica y mostró por separado la revocación actual; una nueva firma
fue rechazada. El retiro del vínculo conservó la evidencia original. Tras
respaldar y restaurar PostgreSQL, invalidar las sesiones anteriores e ingresar
con MFA nueva, se recuperaron los mismos recibos, instantes y bytes del ZIP.
El documento sellado mediante el firmante del servidor mantuvo su clasificación
histórica y su paquete sin cambios.

El corte de interfaz aprobó 872 pruebas Node en 4.169 segundos y seis escenarios
de navegador con HTTP controlado en 19.4 segundos, con un trabajador. Incluyó
revisión del PDF exacto, recuperación explícita, restricciones por rol, retiro
móvil y renovación de la confirmación tras rechazar una firma. La compilación
del workspace y la compilación web aprobaron.

La aceptación posterior de Qadra con servicios reales aprobó dos recorridos
en 20.2 segundos de Playwright, con un trabajador: escritorio de 1440 por
1000 píxeles y móvil de 390 por 844. El navegador descargó el PDF y los bytes
canónicos, recibió las firmas generadas externamente con OpenSSL y confirmó
la operación. La prueba interceptó una respuesta exitosa después del commit
para simular su pérdida. Tras revocar la sesión, observar el rechazo de acceso
e ingresar con MFA nueva, la interfaz recuperó el recibo original mediante
consulta, sin repetir el envío. Descargó el ZIP y, después de retirar el
vínculo, mostró la validez histórica y el retiro posterior por separado.

La primera ejecución detectó una carrera del instrumento: leía la respuesta
capturada antes de terminar la intercepción HTTP. La corrección espera su
finalización conservando la exigencia de respuesta exitosa y recibo exacto.
Los dos recorridos finales aprobaron sin errores de página. Esta aceptación
de navegador no repite la revocación por CRL ni la restauración SQL del
recorrido HTTP anterior. Las ocho capturas se inspeccionaron para comprobar
legibilidad de controles, referencias y resultados en ambos tamaños. La
visualización embebida del PDF no se acredita mediante capturas headless;
los bytes descargados y sus firmas se comprobaron por separado. Los controles
de la revisión publicada y el despliegue conservan su evidencia propia.

La TSA local acredita el resumen del contenido bajo su confianza configurada;
no prueba que las firmas personales existieran al generar el token. Registro
y aceptación conservan los instantes de la aplicación. Las claves sintéticas
permiten comprobar el protocolo, sin demostrar custodia exclusiva de una
persona, identidad civil, FIREL ni una constancia NOM-151 emitida por un PSC.
